We predict the empirical Glicko-2 difficulty rating of a Lichess chess puzzle from its position and solution using a compact, CPU-only pipeline. Existing peer-reviewed work in this area is dominated by either transformer regressors requiring GPU training (GlickFormer, MAE 217.7 on 4.2M puzzles) or shallow rules of thumb (Aimchess), leaving no reproducible gradient-boosting reference with rich feature engineering. We propose \ours, a 155-dimensional feature-engineered histogram gradient-boosted regressor whose novelty over prior CPU-only baselines is a 14-feature \emph{tactical-density} block: per-ply legal-move, legal-capture, and legal-check counts along the solution together with setup-move descriptors. These features directly encode candidate-move ambiguity and forcing structure --- the axis that separates a 1500-rated puzzle from a 2000-rated one. A convex blend of isotonic and CDF quantile-matching calibrators, selected under a validation-MAE budget, restores tail spread while preserving Spearman rank order. On a 50{,}000-puzzle held-out slice, \ours reaches test MAE 240.8 \elo (38\% below the global-mean baseline; 12.8 \elo below a 141-feature parent pipeline), Spearman $\rho = 0.768$, and middle-decile calibration 85.6 \elo, clearing the \textless\,100 stretch target for the first time. Cross-seed reproducibility is essentially exact ($\Delta\text{MAE} = 0.12$), and both seeds independently select the same calibrator blend weight $\alpha = 0.30$. The tactical-density block enters the top-15 correlation-proxy importance list on introduction, closing 12.8 \elo of the residual gap to a much heavier transformer baseline at roughly two orders of magnitude less compute.
